\documentstyle[% 


righttag% 


,11pt% 


,amssymb% 


,russian%               remove in Eng. 


,izvru%                 Set izven% in Eng. 


]{amsart} 


 


% seek & repl {,} for {.}


% verify \sim in (33) 


\newcommand{\re}{\operatorname{Re}} 


\newcommand{\im}{\operatorname{Im}} 


\newcommand{\ch}{\operatorname{ch}} 


\newcommand{\sh}{\operatorname{sh}} 


\newcommand{\tts}[1]{}  


 


 


%\hoffset=-15mm%                  For Eng. 


\setcounter{page}{26} 


\god{2003} 


\nomer{10~(497)} %                        Remove in Eng. 


%\nomer{Vol.\,47, No.\,10}%               Set in Eng. 


%\str{\pageref{begin}--\pageref{end}}%  Set in Eng. 


\udk{517.518} 


 


\author{\it Ф.Н.\,Гарифьянов} 


% NB:   ^^ remove in Eng. 


\title{О биортогональных системах, порождаемых некоторыми  


инволютивными операторами} 


\thanks{Работа выполнена при финансовой поддержке Российского фонда


фундаментальных исследований (код проекта 02-01-00914).} 


 


\date{} 


%\pagestyle{myheadings}%         Set in Eng. 


\pagestyle{plain} %              Remove in Eng. 


\begin{document} 


%\thispagestyle{plain}%          Set in Eng. 


\maketitle 


\label{begin} 


 


 


 


\section*{Введение} 


 


Известно, что оператор сингулярного интегрирования 


$$ 


(S\varphi)(t)\equiv\frac{1}{\pi i}\int_\Gamma(\tau-t)^{-1}


\varphi(\tau)d\tau, 


$$ 


где интеграл понимается в смысле главного значения по Коши, при достаточно 


общих предположениях на плотность $\varphi(t)$ и простую замкнутую кривую 


$\Gamma$ 


является инволютивным оператором. Достаточно потребовать, например, чтобы 


$\varphi(t)\in L_p(\Gamma)$, $p>1$ ([1], с.\,13). Множество неподвижных точек 


оператора 


совпадает с множеством функций, голоморфных внутри $\Gamma$, а решения 


уравнения 


\begin{equation} 


S\varphi=-\varphi 


\end{equation} 


голоморфны вне $\Gamma$ и исчезают на бесконечности. 


 


В данной работе изучаются свойства биортогональных систем, 


порождаемых инволютивными операторами, ядра которых имеют более сложную 


структуру, чем ядро Коши. Пусть $R$ --- квадрат с вершинами в точках 


$t_1=-t_3=-\frac{1}{2}(1+i)$, $t_2=-t_4=\frac{1}{2}(1-i)$, перечисленными в 


порядке 


положительного обхода границы, и $\Gamma=\partial R$. Цель статьи --- 


исследование 


интегрального оператора 


\begin{equation} 


(S\varphi)(t)\equiv\frac{1}{\pi i}\int_\Gamma 


E(\tau-t)\varphi(\tau)d\tau 


\end{equation} 


с нечетной ядерной функцией 


\begin{equation} 


E(z)=z^{-1}+(z+1+i)^{-1}+(z+1-i)^{-1}+(z-1+i)^{-1}+ 


(z-1-i)^{-1}, 


\end{equation} 


а также порождаемых им биортогонально сопряженных систем. Заметим, что 


оператор (2) является сингулярным возмущенным оператором с неподвижными 


особенностями (общую теорию см., напр., в [1], [2]). 


 


В разделе 1 проводится равносильная регуляризация уравнения (1) и 


описывается множество ее решений. В разделе 2 вводится система функций 


\begin{equation} 


\{\varphi_m\}:S\varphi_m=-\varphi_m+q^+_m;\quad 


q_m(z)=(-1)^m(m!)^{-1}(z^m-p_m),\quad m=1,2,\dotsc, 


\end{equation} 


где $p_m$ --- некоторые постоянные, определяемые ниже. Доказывается, что она 


биортогонально сопряжена на $\Gamma$ с системой 


\begin{equation} 


\{E^{(k)}(\tau)\},\quad k=1,2,\dotsc, 


\end{equation} 


в том смысле, что 


\begin{equation} 


(\varphi_m,E^{(k)})\equiv\frac{1}{2\pi i}\int_\Gamma 


\varphi_m(\tau)E^{(k)}(\tau)d\tau=\delta_{m,k}. 


\end{equation} 


В разделе 3 получены биортогональные разложения по системе функций 


\begin{equation} 


\{F_k=E^{(k)}-f^{(k)},\ k=1,2,\dotsc;\ F_0\equiv1\}. 


\end{equation} 


Здесь $f(z)$ --- квазиэллиптическая функция, имеющая в $\ov R$ единственный 


простой 


полюс в нуле с вычетом, равным 1, причем 


\begin{equation} 


f(z+1)=f(z+i)=-f(z). 


\end{equation} 


Изложен также общий принцип трансформации биортогональных систем 


оператором (2), обсужденo возможноe обобщениe полученных 


результатов, указана их связь с теорией целых функций экспоненциального типа 


(ц.\,ф.\,э.\,т) и преобразования Бореля. 


 


В статье приняты следующие обозначения: $l_k$, $k=\ov{1,4}$, --- стороны 


квадрата, пронумерованные в порядке положительного обхода границы, начиная с 


нижней стороны; $\alpha(t)=\{t+i^k$, $t\in l_k\}$ --- обратный сдвиг 


Карлемана ($\alpha(t):\Gamma\to\Gamma$ 


с изменением ориентации), индуцированный порождающими преобразованиями 


соответствующей двоякопериодической группы и разрывный в вершинах; 


$W:\varphi(t)\to\varphi[\alpha(t)]$ --- инволютивный оператор сдвига; 


$(Vf)(z)\equiv f(z+1+i)+f(z-1-i)+f(z-1+i)+f(z+1-i)$ --- линейный 


разностный оператор с постоянными коэффициентами. 


 


 


\section{О решениях уравнения (1)} 


 


Будем считать, что оператор (2) действует на 


линейном пространстве $Q$ --- 


множестве функций, удовлетворяющих условию Г\"ельдера на каждой закрытой 


стороне квадрата, причем сумма скачков в противоположных вершинах равна 


нулю. Такой выбор обусловлен тем, что $S:Q\to Q$. Доказательство этого 


утверждения 


(т.\,е. выяснение характера возможных особенностей у интеграла (2)) 


проводится по схеме, указанной в ([3], с.\,9). Для интеграла $(S\varphi)(t)$ 


справедливы 


формулы Сохоцкого $S^\pm=S\pm J$. Отсюда непосредственно вытекает, что 


неподвижными точками оператора $S$ являются только граничные значения 


аналитических в $R$ функций. 


 


При исследовании уравнения (1) возникают трудности, обусловленные тем, 


что интеграл $(S\varphi)(t)$ вне $R$ является функцией не голоморфной, а 


кусочно-голоморфной с линиями разрывов $\Gamma\pm i\pm1$ (знаки не


согласованы). Поэтому 


применим следующий прием: вместо исходного возьмем неоднородное уравнение 


\begin{equation} 


S\varphi=-\varphi+2g^+, 


\end{equation} 


где $g(z)$ отлична от постоянной и $g^+\in Q$. Другие требования на эту 


функцию 


будут наложены ниже. Очевидно, одним из решений уравнения (9) является 


функция $g^+$. Остается найти еще одно решение $\varphi$, которое заведомо не 


будет 


граничным значением аналитической в $R$ функции. Тогда нетривиальным 


решением уравнения (1) является разность 


\begin{equation} 


\theta=\varphi-g^+. 


\end{equation} 


Разумеется, надо будет доказать, что любое решение (1) представимо в виде 


(10), причем разным функциям $g$ соответствуют разные решения $\theta$. 


 


Построим решения уравнения (9), удовлетворяющие условию 


\begin{equation} 


\varphi=-W\varphi. 


\end{equation} 


Такие функции голоморфны лишь в случае $\varphi\equiv0$. Перепишем с учетом 


(11) 


соотношение (9) в виде $-WSW\varphi=\varphi+2Wg^+$. Вычитая из уравнения (9) 


полученное, имеем 


\begin{equation} 


T\varphi=2g^+-2Wg^+;\quad T\equiv2J+S+WSW. 


\end{equation} 


Ограниченность ядра этого уравнения 


\begin{equation} 


K(t,\tau)=E(\tau-t)-E[\alpha(\tau)-\alpha(t)] 


\end{equation} 


была установлена в [4]. Его частные производные также ограничены на $\Gamma$. 


 


 


\begin{thm} 


Неоднородное уравнение Фредгольма $(12)$ разрешимо, причем 


всегда существует такое его решение, что $\varphi=-W\varphi$, $\varphi\in Q$ 


и 


\begin{equation} 


\int_\Gamma \varphi(t)dt=0. 


\end{equation} 


Фундаментальная система решений $($ф.\,с.\,р.$)$ соответствующего однородного 


уравнения 


\begin{equation} 


T\varphi=0 


\end{equation} 


состоит из единственной нечетной и разрывной в вершинах функции 


$\varphi_0(\tau)$, 


причем $\varphi_0\in Q$, $W\varphi_0=-\varphi_0$, $S(\varphi_0)=-\varphi_0$ и 


\begin{equation} 


\int_\Gamma \varphi_0(\tau)d\tau\ne0. 


\end{equation} 


Ф.\,с.\,р. союзного уравнения 


\begin{equation} 


T'\psi=0 


\end{equation} 


содержит лишь постоянную. 


\end{thm} 


 


 


{\bf Доказательство} разобьем на ряд лемм. Предварительно приведем 


несколько достаточно очевидных утверждений, которыми будем пользоваться 


далее,


 


а) $\varphi\in Q\Leftrightarrow W\varphi\in Q$; 


 


б) $\alpha(i\tau)=-i\alpha(\tau)$, $E(i\tau)=-iE(\tau)$, 


$K(it,i\tau)=-iK(t,\tau)$; 


 


в) $(\varphi,S\psi)=-(\psi,S\varphi)$ (см. обозначение в (6)); 


 


г) если $\varphi\in L_2(\Gamma)$ и выполнено (15), то $\varphi\in Q$. 


 


Пусть $B$ --- множество четных функций со свойством (11) или нечетных 


функций со свойством $\varphi=W\varphi$. Тогда 


\begin{equation} 


b(t)\in B\ \Rightarrow\ \int_{l_j}b(t)dt=0. 


\end{equation} 


Четные функции удовлетворяют одному из условий $\varphi(t)=\pm\varphi(it)$, а 


нечетные --- 


одному из условий $\varphi(t)=\pm i\varphi(it)$. 


 


 


\begin{lem} 


Ф.\,с.\,р. уравнений $(15)$ или $(17)$ не содержат функций из 


множества $B$. 


\end{lem} 


 


 


\begin{pf} 


Допустим, что модуль такого решения принимает свое 


наибольшее значение $M$ в некоторой точке $t\in l_1$. Оценим модуль 


интегрального 


слагаемого в уравнении (15), разбивая его на слагаемые по отдельным сторонам. 


При $\tau\in l_1$ имеем $K(t,\tau)=0$. При $\tau\in l_3$ получим 


$\tau-t=\beta+i$, $\alpha(\tau)-\alpha(t)=\beta-i$, $\beta\in[-1;1]$, 


$K(t,\tau)=-2if(\beta^2)$, 


где $f(x)=(x+1)^{-1}+4(x+5)(x^2+6x+25)^{-1}$. 


Воспользуемся тождеством $\int\limits_{l_3}\varphi(\tau) K(t,\tau)d\tau= 


\int\limits_{l_3}\frac{K(t,\tau)-K(t,-\alpha(\tau))}{2} \varphi(\tau)d\tau$. 


Вычисляя наибольшее значение $|f(\beta_1)-f(\beta_2)|$ при 


$\beta_1,\beta_2\in[0,1]$, придем к 


неравенству $\Big|\int\limits_{l_3}\Big|\le0{,}55M$. 


 


При $\tau\in l_4$ и $\gamma=\tau-t$ имеем 


$$ 


K(t,\tau)=(\gamma+1+i)^{-1}+(\gamma-1-i)^{-1} +(\gamma-1+i)^{-1}


-(\gamma+2)^{-1}-(\gamma+2i)^{-1}-(\gamma+2-2i)^{-1}. 


$$ 


Интегралы с плотностью $\varphi(\tau)$ от соответствующих слагаемых в этой 


сумме 


обозначим через $J_j(t)$. Оценим первый из этих интегралов, ``подправляя'' 


его 


ядро слагаемым, интеграл от которого равен нулю в силу (18), а именно, 


$$ 


|J_1(t)|=\bigg|\int_{l_4}\varphi(\tau)[(\tau-t+1+i)^{-1} 


+(t-1/2-i)^{-1}]d\tau\bigg|= 


\bigg|\int_{l_4}\frac{\tau+1/2}{(\tau-t+1+i)^{-1}(t-1/2-i)} 


\varphi(\tau)d\tau\bigg|\le\frac{M}{3}. 


$$ 


``Подправляющее'' слагаемое выбирается из тех соображений, чтобы появился 


множитель $\tau+0,5$ --- линейный сдвиг, переводящий данную сторону в отрезок 


мнимой оси. Применяя аналогичный прием, получим оценки 


$|J_j(t)|\le\theta_j M$; $\theta_2=1/3$, $\theta_3=\theta_4=\sqrt{5}/5$, 


$\theta_5=\sqrt{3}/3$, $\theta_6=\sqrt{26}/26$. 


Поэтому 


справедливо суммарное неравенство 


$\Big|\int\limits_{l_4}\Big|\le2{,}34{M}$. Точно такая же оценка 


верна и для модуля интеграла по стороне $l_2$. В итоге из уравнения (15) 


имеем 


$2\pi M\le5,23M$ и $M=0$. 


 


Заметим, что четная функция $\varphi\in Q$ непрерывна. Поэтому вместе с четным 


решением $\varphi(t)$ ф.\,с.\,р. содержит и нечетную функцию $\varphi'(t)$. 


Достаточно 


продифференцировать уравнение (15) и в интегральном слагаемом применить 


формулу интегрирования по частям, воспользовавшись равенством 


$\partial K/\partial t=-\partial K/\partial\tau$. Сумма скачков 


произведения $\varphi(\tau)E[\alpha(\tau)-\alpha(t)]$ по 


переменной $\tau$ в противоположных вершинах равна нулю и внеинтегральное 


слагаемое исчезает. 


\end{pf} 


 


 


\begin{lem} 


Ф.\,с.\,р. уравнения $T\varphi=0$ не содержит решений со свойством 


$\varphi=W\varphi$. Ф.\,с.\,р. союзного уравнения содержит единственную такую 


функцию --- постоянную. 


\end{lem} 


 


 


{\bf Доказательство.} В силу леммы 1 такое решение обязательно четно, т.\,е. 


его производная $\varphi'(t)\equiv0$ как решение однородного уравнения из 


множества $B$, а 


постоянная удовлетворяет союзному уравнению. 


 


 


\begin{lem} 


Ф.\,с.\,р. союзного уравнения не содержит функций со свойством 


$\varphi=-W\varphi$. Ф.\,с.\,р. уравнения $T\varphi=0$ содержит единственную 


такую функцию. 


\end{lem} 


 


 


\begin{pf} 


По альтернативе Фредгольма и в силу леммы 2 ф.\,с.\,р. 


уравнения (15) содержит хотя бы одну такую функцию $\varphi_0(\tau)$, а по 


лемме 1 эта 


функция нечетна. Поскольку для такой функции $T=S^+-WS^+$, то 


\begin{equation} 


(S\varphi_0)(z)=0,\quad z\in R. 


\end{equation} 


 


Предположим, что для решения $\varphi_0$ не выполнено условие (16). 


Рассмотрим квазипериодический аналог интеграла типа Коши, введенный в [5]: 


\begin{equation} 


\psi(z)=\frac{1}{\pi i}\int_\Gamma \varphi_0(\tau)\zeta(\tau-z)d\tau, 


\end{equation} 


где дзета-функция Вейерштрасса построена по примитивным периодам 1 и $i$. В 


силу сделанного допущения интеграл (20) является уже двоякопериодическим 


аналогом интеграла типа Коши. Справедлив аналог формулы Сохоцкого 


$\psi^+=\varphi_0-W\varphi_0+\psi$, т.\,е. в данном случае 


$\psi^+=2\varphi_0+\psi$. Имеем 


$(S\psi)(z)=(S\psi^+)(z)=2\psi(z)$, $z\in R$ (см. (19)). Складывая равенства 


$S^+\psi=2\psi^+$ и 


$WS^+\psi=2W\psi^+$, получим $S^+\psi+WS^+\psi=4\psi$ $\Rightarrow$ 


$T'\psi=0$, т.\,е. приходим к 


противоречию с леммой 1. Итак, решение $\varphi_0$ удовлетворяет (16), т.\,е. 


единственно. 


 


Заметим, что $\varphi_0(i\tau)=-i\varphi_0(\tau)$ и $\varphi_0\notin 


C(\Gamma)$. Первое утверждение 


выполнено в силу (16). Если предположить, что данная функция непрерывна, то 


ее производная была бы четным решением из класса $B$ и по лемме 1 имели бы 


$\varphi'_0\equiv0$. 


\end{pf} 


 


 


Получим формулу для вычисления решения $\varphi_0$, которое играет в 


дальнейшем важную роль. Введем разность 


\begin{equation} 


\beta(t,\tau)=\zeta(\tau-t)-\zeta[\alpha(\tau)-\alpha(t)]. 


\end{equation} 


Здесь и всюду в дальнейшем, если только это особо не оговорено, уже считаем, 


что дзета-функция Вейерштрасса построена по примитивным периодам $1\pm i$, 


т.\,е. 


по диагоналям квадрата. Разность (21) в зависимости от взаимного расположения 


точек $\tau$ и $t$ на сторонах квадрата принимает девять различных значений: 


$0$, $\pm\eta_1$, $\pm\eta_2$, $\pm(\eta_1\pm\eta_2)$, где 


$\eta_{1,2}=2\zeta\big(\frac{1\pm i}{2}\big)$, т.\,е. является кусочной 


постоянной с точками разрыва в вершинах. Введем еще одну кусочную 


постоянную 


\begin{equation} 


c_\varphi(t)=\frac{1}{\pi i}\int_\Gamma \beta(t,\tau)\varphi(\tau)d\tau. 


\end{equation} 


Если выполнено (14), то кусочная постоянная (22) является просто постоянной 


(непосредственным подсчетом скачков в вершинах убеждаемся, что они равны 


нулю). 


 


Рассмотрим компактный оператор 


\begin{equation} 


(A\varphi)(t)\equiv\frac{1}{\pi i}\int_\Gamma 


D(\tau-t)\varphi(\tau)d\tau 


\end{equation} 


с ядерной функцией 


\begin{equation} 


D(z)=\zeta(z)-E(z), 


\end{equation} 


голоморфной в круге $|z|<2$. Решение $\varphi_0$ удовлетворяет неоднородному 


уравнению $(T_1\varphi)=-c_{\varphi_0}(t)$, $T_1=2J-A+WA$. Оператор $T_1$ 


обладает важным 


преимуществом по сравнению с оператором $T$. 


 


 


\begin{thm} 


Оператор $T_1$ обратим. 


\end{thm} 


 


 


{\bf Доказательство.} Любое решение однородного уравнения 


\begin{equation} 


T_1\varphi=0 


\end{equation} 


удовлетворяет условию (11) из-за структуры его ядра. Непосредственным 


интегрированием уравнения (25) убеждаемся, что выполнено условие (14). 


Кусочная постоянная (22) является в таком случае постоянной, антиинвариантной 


относительно сдвига, т.\,е. просто нулeм. Тогда выполнено равенство (15) и 


$\varphi\equiv0$ 


в силу теоремы 1. 


 


 


\begin{cor} 


Справедливо представление $\varphi_0(t)=\lambda T^{-1}_1a_t$, где кусочная 


постоянная $a_t=\{(-i)^k$, $t\in l_k$; $k=\ov{1,4}\}$, а постоянная $\lambda$ 


выбрана так, что $(\varphi_0,1)=1$. 


\end{cor} 


 


 


Рассмотрим вопрос об эквивалентности уравнений (9) и (12). Было 


доказано, что если\linebreak


$\varphi$ --- решение (9) со свойством (11), то такое решение 


удовлетворяет и уравнению (12). Обратное, вообще говоря, неверно, хотя 


уравнение (12) разрешимо --- ф.\,с.\,р. союзного уравнения (15) содержит 


только 


постоянную. Существует решение, удовлетворяющее условиям (11) и (14), но оно 


может не удовлетворять (9). Действительно, перепишем (12) в виде 


$S^+\varphi-WS^+\varphi=g^+-Wg^+$, откуда 


\begin{equation} 


(S\varphi)(t)=-\varphi(t)+2g^+(t)+c. 


\end{equation} 


Остается выяснить, когда постоянная является нулем, и вопрос о равносильности 


уравнений (9) и (12) будет решен. Умножим обе части уравнения (26) на 


$\varphi_0$ и 


проинтегрируем по $\Gamma$, воспользовавшись равенствами 


$(\varphi_0,S\varphi)=0$ и $(\varphi_0,\varphi)=0$. 


Первое равенство справедливо в силу (19), второе очевидно. Другими словами, 


постоянная исчезает, если 


\begin{equation} 


(g^+,\varphi_0)=0. 


\end{equation} 


Условие (27) впредь считаем выполненным. 


 


Приступим к непосредственному изучению исходного уравнения (1), 


которое перепишем в виде 


\begin{equation} 


S\theta=-\theta. 


\end{equation} 


Напомним, что под функцией $\varphi$ понимается вполне определенное решение 


уравнения (12), удовлетворяющее условиям (11) и (14). 


 


 


\begin{thm} 


Любое решение уравнения $(28)$ представимо в виде $(10)$. 


Разным функциям $g^+$ соответствуют разные решения $\theta$. 


\end{thm} 


 


 


\begin{pf} 


Положим $\varphi=\theta-W\theta$. Без ограничения общности 


считаем, что выполнено (14), т.\,к. функция $\varphi$ определена с точностью 


до слагаемого $\lambda\varphi_0$ и функции $\varphi_0$ в разности (10) 


соответствует $g^+\equiv0$. Тогда при 


$z\in R$ в силу (28) справедливо тождество $S\varphi\equiv-SW\theta$, т.\,е. 


функция $\varphi$ 


удовлетворяет уравнению (12) при $g(z)=-(SW\theta)(z)+\lambda$, где 


постоянная выбрана 


так, чтобы выполнялось равенство (27). 


 


Теперь предположим, что конечной системе линейно независимых 


функций $\{g^+_j\}$ соответствует линейно зависимая система решений 


$\{\theta_j\}$. Равенство 


$\sum\lambda_j(\varphi_j-g^+_j)=0$ возможно лишь в случае, когда для функции 


$F^+=\sum\lambda_j g^+_j$ 


имеем $F^+=-WF^+$, т.\,е. $F^+\equiv0$. 


\end{pf} 


 


 


Теорема 3 позволяет получить ряд важных следствий. Обозначим через $Q_1$ 


множество решений уравнения (28), удовлетворяющих условию (14). Множество 


функций $g(z)$, голоморфных в $R$ и удовлетворяющих условиям $g^+\in Q$ и 


(27), 


обозначим через $Q_2$. Ясно, что $Q_1$ и $Q_2$ --- подпространства линейного 


пространства $Q$. Формула (10) определяет некоторый гомеоморфизм $B:Q_2\to 


Q_1$. 


Найдем представление для оператора $B$. Для этого перейдем к описанию 


множества $Q_1$. Уравнение (25) имеет лишь тривиальное решение. Союзное к 


нему 


уравнение с учетом теоремы Коши запишется в виде $T_2\mu=0$, $T_2\equiv 


2J+AW$. 


 


 


\begin{thm} 


Пусть $\varphi$ --- решение неоднородного уравнения $(12)$, 


удовлетворяющее условиям $(11)$ и $(14)$, а $\mu^+$ --- решение уравнения 


\begin{equation} 


T_2\mu^+=2g^+. 


\end{equation} 


Тогда $\varphi=\mu^+-W\mu^+$. 


\end{thm} 


 


 


{\bf Доказательство.} По альтернативе Фредгольма уравнение (29) безусловно 


разрешимо. Запишем его в равносильном виде $2\mu^+-A\varphi=2g^+$, поскольку 


по 


теореме Коши $A\mu^+=0$. Применяя к обеим частям уравнения оператор сдвига и 


вычитая из него полученное, имеем $T_1\varphi=2(g^+-Wg^+)$. Отсюда уже легко 


перейти к уравнению (12), поскольку кусочная постоянная (22) 


исчезает (см. доказательство теоремы 2). 


 


 


\begin{cor} 


Справедливы соотношения $\theta=\mu^+-W\mu^+-g^+$ $\Leftrightarrow$ 


$\theta=Bg^+$, 


где $\mu^+$ удовлетворяет уравнению (29), а $B=-J+T^{-1}_2-WT^{-1}_2$. 


\end{cor} 


 


 


\section{Построение биортогональных систем и некоторые их свойства} 


 


 


Приступим к построению биортогонально сопряженных систем. Выберем 


в формуле (4) постоянные $p_m$ так, чтобы $q^+_m\in Q_2$, т.\,е. выполнялись 


условия (27). 


При $m=0$ положим $p_0=-1$. Система функций (4) со свойством (11) выбрана 


так, чтобы выполнялось и условие (14), $m\ge1$. Систему производных (5) 


пополним единицей и положим 


\begin{equation} 


E_k=E^{(k)},\quad k\ge1;\quad E_0=1. 


\end{equation} 


Систему функций (4) пополним функцией $\varphi_0$. 


 


 


\begin{thm} 


Системы функций $(30)$ и $(4)$ биортогональны на $\Gamma$ в смысле $(6)$. 


\end{thm} 


 


 


\begin{pf} 


При $m=0$ имеем $(\varphi_0,1)=1$, $(\varphi_0,E_j)=0$, $j\ge1$. Первое 


из этих равенств выполняется за счет выбранной нормировки, а остальные --- в 


силу равенства (19). Кроме того, из (14) получим $(\varphi_m,1)=0$, $m\ge1$. 


Во всех 


остальных случаях имеем $(S\varphi_m)(z)=2q_m(z)$, $z\in R$, $\Rightarrow$ 


$(\varphi_m,E_j)=\delta_{m,j}$. Механизм 


образования биортогональности совершенно элементарен и основан на 


простейших равенствах $q^{(k)}_m(0)=\{0$, $k\ne m;m!$, $k=m\}$, $k\ge1$. 


\end{pf} 


 


 


\begin{thm} 


Система функций $(4)$ полна в том смысле, что если $g^+\in Q_2$ и 


$(g^+,\varphi_m)=0$, $m\ge1$, то $g^+\equiv0$. 


\end{thm} 


 


 


\begin{pf} 


Предположим противное. Обозначим через $\varphi$ решение 


соответствующего уравнения (9), удовлетворяющее условиям (11) и (14). 


Умножим обе части уравнения (9) на функции (4) и проинтегрируем по $\Gamma$. 


В 


итоге получим $(\varphi,\tau^m)=0$ $\Rightarrow$ $\varphi=\varphi^+$, а с 


учетом (11) имеем $\varphi\equiv0$ и $g\equiv0$. 


\end{pf} 


 


 


\begin{note} 


Теорема останется справедливой, если функции (4) заменить 


на собственные функции $\theta_m=Bq_m$ оператора (2). 


\end{note} 


 


 


\begin{thm} 


Пусть функция $b(z)$ голоморфна вне $\Gamma$, исчезает на 


бесконечности и $b^-\in Q$. Тогда система $\{\theta_m\}$ полна в том смысле, 


что если 


$(b^-,\theta_m)=0$, $m\ge1$, то $b^-\equiv0$. 


\end{thm} 


 


 


\begin{pf} 


С учетом соотношения (28) перепишем условие теоремы 


в виде $(b^-,S\theta_m)=-(\theta_m,Sb^-)= -(\theta_m,Sb^-+b^-)= 


-(\theta_m,Vb)$, причем последнее 


равенство вытекает из теоремы Коши и функция $(Vb)(z)$ голоморфна в $R$. 


Если 


ее подправить некоторой постоянной (для функции исходной системы выполнено 


(14)), то она будет принадлежать подпространству $Q_2$. Применяя теорему 6 и 


замечание 1, придем к линейному разностному уравнению 


\begin{equation} 


(Vb)(z)=\lambda,\quad z\in R. 


\end{equation} 


Сумма в левой части уравнения (31) голоморфна в связной внешности четырех 


конгруэнтных квадратов без общих точек,


которая содержит наряду с $R$ бесконечно 


удаленную точку. Поэтому постоянная исчезает и $b(z)\equiv0$ ([6], с.\,229). 


\end{pf} 


 


 


\begin{cor} 


Если $\psi^-\in Q_1$ и $\psi(\infty)=0$, то $\psi\equiv0$. 


\end{cor} 


 


 


Рассмотрим интегралы типа Коши 


\begin{equation} 


\Omega_m(z)=\frac{1}{2\pi i}\int_\Gamma(\tau-z)^{-1} 


\theta_m(\tau)d\tau, 


\end{equation} 


где $\theta_0=\varphi_0$, $\theta_m=Bq_m$, $m\ge1$ (см. замечание 1 и 


следствие 2). Это 


кусочно-голоморфные функции с линией разрывов $\Gamma$, причем 


$\Omega^+_m-\Omega^-_m=\theta_m$. Далее 


используем пространства аналитических функций $A(D)$, $A(cD)$, $A[D]$, 


$A[cD]$ [7]. 


 


 


\begin{thm} 


Если область $D\subset R$, то система $(32)$ полна в $A(D)$. 


\end{thm} 


 


 


\begin{thm} 


Если $R\subset D$, то система $(32)$ полна в $A(cD)$. 


\end{thm} 


 


 


{\bf Доказательство} непосредственно следует из теорем 5 и 6 в силу 


известного критерия Банаха. Подчеркнем еще раз, что речь идет о разных 


системах функций $\{\Omega^+_m\}$ и $\{\Omega^-_m\}$, связанных на $\Gamma$ 


формулой Сохоцкого. 


 


Элементы пространства $Q_1$ из-за нечетности ядерной функции 


ортогональны в том смысле, что $(f,\psi)=0$ $\forall f,\psi\in Q_1$. 


Справедлив в 


некотором смысле обратный результат: $f\in Q_1$ $\Leftrightarrow$ 


$(f,\psi)=0$ $\forall \psi\in Q_1$ $\Leftrightarrow$


$(f,\theta_m)=0$ $\forall m$. 


Действительно, из (28) и последнего равенства имеем 


$(f,S\theta_m)=0$ $\Rightarrow$ $(\theta_m,Sf)=0$, т.\,е. для функции 


$g^+=S^+f$ выполнены условия 


теоремы 6 и $f\in Q_1$. 


 


Пусть $f$ --- квазиэллиптическая функция с простым полюсом в нуле 


(вычет равен 1), удовлетворяющая условиям (8), т.\,е. $Wf=-f$. Систему 


мероморфных в $R$ (за исключением $E_0$) функций (30) заменим системой 


функций (7), которая также биортогональна на $\Gamma$ системе решений (4) в 


силу своей структуры. Она состоит из голоморфных в $\ov R$ функций, чем 


выгодно 


отличается от функций системы (30). Действительно, у функций (7) полюсы


расположены в 


точках $mi+n$, $m^2+n^2\ne0$, и 2. 


 


 


\section{Свойства биортогональных рядов. Приложения} 


 


 


Обозначим через $D_0$ область, ограниченную дугами окружностей 


$|z-i^k|=\sqrt{2}/2$, $k=\ov{1,4}$, и содержащую начало координат. Возьмем 


четную 


функцию $\Phi\in A[cD_0]$ и поставим ей в соответствие биортогональный ряд 


\begin{equation} 


\Phi\sim\sum^\infty_{m=0}\alpha_m\Omega_m 


\end{equation} 


с естественным образом определенными коэффициентами $\alpha_m=(\Phi,F_m)$. 


 


 


\begin{thm} 


Ряд $(33)$ сходится к функции, по которой он построен, 


равномерно на замыкании $cD_0$. 


\end{thm} 


 


 


{\bf Доказательство} идентично рассуждениям, приведенным в [8]. 


Разность между $\Phi$ и суммой биортогонального ряда (33) на $\Gamma$ 


является четной 


г\"ельдеровской функцией $\psi$. Поскольку $\forall j$ \ $(\psi,F_j)=0$ в 


силу равномерной 


сходимости ряда на $\Gamma$, то здесь вместо $\psi$ можно взять функцию 


$\varphi=\psi+a^+$ и за 


счет подбора голоморфной в $R$ функции $a(z)$ потребовать, чтобы 


$\varphi=-W\varphi$ $\Rightarrow$ $T\varphi=0$. Тогда в силу (14) имеем 


$\varphi\equiv0$ и по теореме Лиувилля 


$\psi\equiv0$. 


 


 


\begin{note} 


Ряд (33) на $cD_0$ сходится абсолютно. 


\end{note} 


 


 


\begin{note} 


В случае произвольной функции $\Phi\in A[cD_0]$ теорема 10 и 


замечание 2 остаются справедливыми, если замыкание $cD_0$ заменить на 


произвольный компакт $\ov H\subset\ov{cD_0}$. 


\end{note} 


 


 


Теперь возьмем четную функцию $a(z)\in A(D_0)$, $a^+\in C^2(\partial D_0)$ и 


поставим ей в соответствие ряд 


\begin{equation} 


a\sim\sum^\infty_{m=0}\beta_m F_m 


\end{equation} 


с естественным образом определенными коэффициентами 


$\beta_m=(a^+,\Omega^-_m)$. 


 


 


\begin{thm} 


Ряд $(34)$ сходится к функции, по которой он построен, 


равномерно на замыкании $\ov D_0$. 


\end{thm} 


 


 


{\bf Доказательство} с учетом предыдущей теоремы проводится по 


стандартной схеме, основанной на разложении ядра Коши в биортогональный ряд 


([9], гл.\,4, \S\,6, п.\,3). Справедливы также замечания, аналогичные 


замечаниям к теореме 10. 


 


Теоремы 4 и 5 позволяют ввести системы 


\begin{equation} 


\{WF_m\},\quad \{\mu^+_k\}, 


\end{equation} 


где $\mu^+_k=2T^{-1}_2q_k$, $m,k\ge1$. Они также биортогональны на $\Gamma$ в 


смысле (6). 


Заметим, что $\mu^+_k=\mu_k$, т.\,к. эти функции голоморфны в круге 


$|z|<1,5$, 


содержащем $\Gamma$. Это прямо следует из вида ядерной функции (24) и формулы 


(29). 


Заменим (с сохранением биортогональности) первую из систем (35) на систему 


интегралов типа Коши 


$$ 


\bigg\{\frac{1}{2\pi i}\int_\Gamma F_m[\alpha(\tau)] 


(z-\tau)^{-1}d\tau\bigg\},\quad |z|>\sqrt{2}/2. 


$$ 


Полученную систему пополним $\Omega_0$, а вторую --- единицей. 


Кусочно-линейный 


сдвиг аналитически продолжим с соответствующей стороны квадрата в четыре 


сегмента, дополняющих $R$ до круга $|z|\le\sqrt{2}/2$. Для таких 


биортогонально 


сопряженных на окружности $|z|=\sqrt{2}/2$ систем остаются справедливыми 


теоремы 


8 и 9 с заменой области $D_0$ на этот круг. 


 


Вообще, рассмотрим трансформации оператором $S$ произвольных 


биортогонально сопряженных на $\Gamma$ систем $\{a_m\}$, $\{b_k\}$; $a_m\in 


A(R)$, $b_k\in A(cR)$. 


Предполагаем, что для функций второй системы выполнено условие (14), т.\,е. 


$b^-_k=\gamma^+_k+W\sigma^+_k$, где $\gamma,\sigma\in A(R)$. Из теории 


краевых задач со сдвигом Карлемана и 


г\"ельдеровскими коэффициентами следует, что такое представление существует и 


единственно. Вторую систему заменим системой $\{W\sigma^+_k\}$ с сохранением 


биортогональности. Тогда непосредственно проверяется, что система 


антиинвариантных относительно сдвига функций 


\begin{equation} 


\varphi_m:S\varphi_m=-\varphi_m+a^+_m 


\end{equation} 


биортогональна системе $\{2S^+\sigma^+_k-2S^+W\sigma^+_k\}$. Заметим, что 


функции первой 


системы в силу (14) определяются с точностью до произвольной постоянной, за 


счет выбора которой добьемся разрешимости уравнения (36). 


 


Пусть теперь $a(z)$ и $b(z)$ --- голоморфные соответственно внутри и вне 


$\Gamma$ 


функции, совпадающие на этих связных компонентах с интегралом типа Коши 


$\Omega_0$ (см. (32)). Тогда уравнение (1) запишется в виде 


\begin{equation} 


a(z)=-(Vb)(z),\quad z\in R. 


\end{equation} 


Соотношение (37) несколько неудобно тем, что оно связывает значения двух 


функций с разными областями голоморфности (правда, представимые в виде 


одного интеграла типа Коши с плотностью, антиинвариантной относительно 


сдвига). Справедлива формула Сохоцкого $a^+-b^+=\varphi_0$. Исключая отсюда 


плотность и используя соотношение (37), придем к частному случаю задачи 


\begin{equation} 


(Mb^-)(t)+G(t)(WMb^-)(t)=g(t),\quad t\in\Gamma, 


\end{equation} 


при $G\equiv1$ и $g\equiv0$. Здесь $M=J+V$. 


Задачу (38) трактуем как внешнюю 


многоэлементную краевую задачу, содержащую наряду со сдвигом Карлемана 


еще и сдвиги, переводящие контур во внутренность области. Решение ищется в 


классе функций, голоморфных вне $\Gamma$ и исчезающих на бесконечности, 


граничные 


значения которых на каждой открытой стороне квадрата удовлетворяют условию 


Г\"ельдера, а в вершинах допускаются, самое большее, логарифмические 


особенности. Коэффициенты задачи $G$ и $g$ удовлетворяют всем ограничениям, 


накладываемым на них при постановке классической задачи Карлемана со 


сдвигом ([10], гл.\,IV, \S\,13, п.\,1), 


в том числе и условиям, гарантирующим от 


переопределенности задачи. Равносильную регуляризацию задачи (38) в общем 


случае можно провести по схеме, предложенной в [4], если только 


$G(t_j)=1$, $j=\ov{1,4}$. При этом существенно используется компактность 


оператора 


(23). Более подробно остановимся на частном случае $G\equiv1$ (т.\,е. $g=Wg$), 


поскольку в этом случае исследование задачи тесно связано со свойствами ранее 


введенных операторов. Будем искать решение в виде 


\begin{equation} 


b(z)=\lambda z^{-1}+\frac{1}{\pi i}\int_\Gamma \psi(\tau) 


(\tau-z)^{-1}d\tau,\quad z\in c\ov R, 


\end{equation} 


где $\lambda$ --- некоторая постоянная, а $\psi=W\psi$ --- неизвестная 


плотность. Структура 


представления (39) обусловлена тем, что интеграл типа Коши с такой плотностью 


имеет на бесконечности нуль не менее чем второго порядка. Любое решение 


указанного частного случая задачи (38) представимо в виде (39). Это прямо 


следует из того факта, что задача Карлемана $F^+-WF^+=\varphi$ имеет 


единственное 


условие разрешимости (14). В итоге вместо задачи (38) имеем интегральное 


уравнение Фредгольма второго рода 


\begin{equation} 


T'\psi=g_1, 


\end{equation} 


где $g_1(t)=\lambda t^{-1}+\lambda\alpha(t)^{-1}-g(t)$. Из ранее 


установленных свойств функции $\varphi_0$ 


следует, что уравнение (40) разрешимо тогда и только тогда, когда 


$(g_1,\varphi_0)=0$ $\Leftrightarrow$ $(g,\varphi_0)=0$ $\Leftrightarrow$ 


$(g^+,\varphi_0)=0$. 


Здесь $g^+$ определяется из безусловно разрешимой 


краевой задачи Карлемана $g^++Wg^+=g_1$, решение которой можно выписать в 


явном виде ([5], с.\,26). Если уравнение (40) разрешимо, то любое его решение 


обладает свойством $\psi=W\psi$ и определено с точностью до произвольной 


постоянной. Оператор $T'$ несколько неудобен тем, что он не является 


обратимым. 


Поэтому, рассуждая аналогично доказательству теоремы 4, можно получить 


следующий результат: пусть $\mu^+=T^{-1}_3g^+$, где $T_3=2J-AW$, тогда 


$\psi=\mu^++W\mu^++c$. 


 


Вернемся к исходному соотношению (37). Поскольку функция $b(z)$ 


голоморфна вне $R$ и исчезает на бесконечности, то функция $a(z)$ голоморфна 


вне 


четырех конгруэнтных квадратов, имеющих общие вершины с $R$, причем 


$a(\infty)=0$. Поэтому и функция $a(z)$, и функция $b(z)$ являются нижними 


функциями, ассоциированными по Борелю 


соответственно с ц.\,ф.\,э.\,т. $A(z)$ и $B(z)$. 


Сопряженной индикаторной диаграммой (наименьшим выпуклым множеством, 


содержащим все особенности нижней функции) для первой функции является 


квадрат $3R$, для второй --- квадрат $R$. В силу ранее установленных свойств 


функции $\varphi_0$ искомые функции представимы в виде $A_1(z^4)$ и 


$B_1(z^4)$, где $A_1(z)$ и 


$B_1(z)$ --- целые функции порядка $\rho=0,25$ и типа $\sigma=\sqrt{2}/2$. 


Умножая обе части 


уравнения (37) на множитель $\exp(\lambda z)dz$ и интегрируя по окружности 


$|z|=r$ 


достаточно большого радиуса (т.\,е. применяя преобразование Бореля в 


окрестности бесконечно удаленной точки), имеем 


\begin{equation} 


A(z)=4\cos z\ch z\,B(z). 


\end{equation} 


Индикаторы искомых целых функций являются $\pi/2$-периодическими 


функциями и полностью определяются из формул 


$$ 


h_B(\theta)=\frac{1}{2}(\cos\theta+\sin\theta),\quad 


\theta\in[0,\pi/2];\quad h_A(\theta)=3h_B(\theta). 


$$ 


 


 


\begin{thm} 


Целые функции $A(z)$ и $B(z)$ являются функциями вполне 


регулярного роста $($в.\,р.\,р.$)$ во всей плоскости, 


т.\,е. множество их корней 


правильно распределено {\em ([11], гл.\,II--III)}. 


\end{thm} 


 


 


\begin{pf} 


С учетом формулы (41) достаточно ограничиться 


рассмотрением целой функции 


\begin{equation} 


B(z)=-\int_\Gamma \varphi_0(\tau)\exp(z\tau)d\tau. 


\end{equation} 


Дважды применяя формулу интегрирования по частям, имеем 


$$ 


B(z)=-z^{-1}\varphi_0(\tau)\exp(z\tau)\big|_\Gamma 


+z^{-2}\bigg[\varphi'_0(\tau)\exp(z\tau)\big|_\Gamma 


-\int_\Gamma \varphi''_0(\tau)\exp(z\tau)d\tau\bigg]. 


$$ 


Луч $\arg z=\pi/4$ является одним из четырех лучей, на которых значение 


индикатора равно типу (направление наибольшего роста). Первое из приращений 


не исчезает из-за разрывности плотности в вершинах. Именно оно определяет 


асимптотическое поведение целой функции на данном луче, поскольку любая 


производная плотности ограничена в силу структуры ядра (13). Итак, функция 


(42) в.\,р.\,р. на данном луче. 


Поскольку ее индикатор является тригонометрическим в 


любом из координатных углов, то она в.\,р.\,р. 


внутри этих углов ([11], гл.\,III, \S\,7, 


теорема~6). Кроме того, множество лучей в.\,р.\,р. 


целой функции либо пусто, либо замкнуто ([11], гл.\,III, \S\,1), 


т.е. функция (42) в.\,р.\,р. во всей плоскости. 


\end{pf} 


 


 


\begin{cor} 


Множество корней ц.\,ф.\,э.\,т. (42) внутри каждого из 


координатных углов имеет нулевую плотность ([11], с.\,202). Другими словами, 


``почти все'' корни целых функций (41) и (42) лежат в ``окрестностях'' 


координатных осей. 


\end{cor} 


 


 


Теперь воспользуемся одной из формул преобразования Бореля ([12], \S\,1, 


теорема 1.1.4) 


\begin{equation} 


b(z)=\int^\infty_0 B(x)\exp(-z x)dx,\quad \re z>0,5. 


\end{equation} 


В уравнении (37) для двух слагаемых $b(z+1\pm i)$ эта формула справедлива, 


т.\,к. 


$z\in R$. Для двух других слагаемых $b(z-1\pm i)$ формула неприменима, т.\,к. 


несобственный интеграл (43) расходится. Однако здесь можно воспользоваться 


нечетностью функции $b(z)$ и переписать уравнение (37) в виде 


$$ 


a(z)=4\int^\infty_0 B(x)\exp(-x)\cos x\sh(z x)dx,\quad 


|\re z|<0,5. 


$$ 


Заменяя здесь $z$ на $iz$ с учетом условия $a(iz)=-ia(z)$ и пользуясь 


формулой 


обращения синус-преобразования Фурье, имеем 


$$ 


\pi A(x)=-2\cos x\ch x\int^\infty_0 


a(t)\sin(x t)dt. 


$$ 


Последнее равенство справедливо лишь на вещественной оси, где интеграл 


сходится абсолютно и равномерно, т.\,к. нижняя функция имеет на бесконечности 


нуль третьего порядка и голоморфна в полосе $|\im z|<0,5$. 


 


В заключение рассмотрим вопрос об обобщении полученных результатов. 


Ядерная функция (3) представляет собой сумму первых пяти слагаемых 


разложения дзета-функции Вейерштрасса, построенной по примитивным 


периодам $1\pm i$, в ряд простейших дробей. Если в формуле (3) взять большее 


число 


слагаемых (среди которых обязательно есть пять уже указанных) так, чтобы 


опять 


выполнялось условие $E(iz)=-iE(z)$, то создается впечатление, что основные 


результаты статьи можно перенести на этот случай. Однако это, вообще говоря, 


не так. Дело в том, что при доказательстве леммы 1, которая обеспечивала 


``хорошую разрешимость'' уравнения (12), применялись специфические оценки, 


существенно использующие конкретный вид ядерной функции (3). Все же можно 


утверждать следующее: если таких слагаемых ``достаточно много'' (но конечное 


число), то картина разрешимости уравнения (12) по-прежнему описывается 


теоремой 1 (см. по этому поводу схему рассуждений, впервые предложенную в 


[13]). При этом возможны два случая. Рассмотрим внешность всех квадратов, в 


которые преобразования группы, вошедшие в ядерную функцию, переводят $R$. 


Если это множество связно или распадается на связные компоненты, одна из 


которых вместе с $R$ содержит и бесконечно удаленную точку, то все результаты 


статьи и схема, по которой они получены, переносятся на данный случай (с 


понятными изменениями формул (41)--(43) и т.\,д.). 


В других случаях теоремы 7 и 


9, по всей видимости, неверны, а рассуждения, проведенные после формулы (37), 


вообще не имеют смысла, т.\,к. наличие такой связной компоненты при выводе 


этих результатов имело определяющее значение. 


 


Автор выражает благодарность проф. Л.А.\,Аксентьеву за ряд ценных 


советов и замечаний. 


 


 


\begin{thebibliography}{99} 


 


\bibitem{1} 


Карапетянц Н.К., Самко С.Г. {\em Уравнения с инволютивными операторами и их 


приложение}. -- Ростов-на-Дону: Изд-во Ростовск. ун-та, 1988. -- 192 с. 


\bibitem{2} 


Дудучава Р.В. {\em Интегральные уравнения свертки с разрывными 


предсимволами, сингулярные интегральные уравнения с неподвижными 


особенностями и их приложения к задачам механики} // Тр. Тбил. матем. ин-та 


АН ГрузССР. -- 1974. -- Т.\,60. -- 134 с. 


\bibitem{3} 


Гарифьянов Ф.Н. {\em Проблема обращения особого интеграла и разностные 


уравнения для функций, аналитических вне квадрата} // Изв. вузов. 


Математика. -- 1993. -- \No\,7. -- С.\,7--16. 


\bibitem{4} 


Гарифьянов Ф.Н. {\em Интегральное представление аналитической функции 


внутри параллелограмма и его приложение} // Изв. вузов. Математика. -- 1991. 


-- \No\,12. -- С.\,8--12. 


\bibitem{5} 


Чибрикова Л.И. {\em О граничных задачах для прямоугольника} // Учен. зап. 


Казанск. ун-та. -- 1963. -- Т.\,123. -- \No\,9. -- С.\,15--39. 


\bibitem{6} 


Коробейник Ю.Ф. {\em О решениях некоторых функциональных уравнений в 


классах функций, аналитических в выпуклых областях} // Матем. сб. -- 1968. -- 


Т.\,75. -- \No\,2. -- С.\,225--234. 


\bibitem{7} 


Хавин В.П. {\em Пространства аналитических функций} // Итоги науки и 


техн. Матем. анализ. -- М.: ВИНИТИ, 1966. -- С.\,76--174. 


\bibitem{8} 


Гарифьянов Ф.Н. {\em Трансформации биортогональных систем и некоторые их 


приложения}. I // Изв. вузов. Математика. -- 1996. -- \No\,6. -- С.\,5--16. 


\bibitem{9} 


Леонтьев А.Ф. {\em Ряды экспонент}. -- М.: Наука, 1976. -- 536 с. 


\bibitem{10} 


Литвинчук Г.С. {\em Краевые задачи и сингулярные интегральные уравнения со 


сдвигом}. -- М.: Наука, 1977. -- 448 с. 


\bibitem{11} 


Левин Б.Я. {\em Распределение корней целых функций}. -- М.: ГИТТЛ, 


1956. -- 632 с. 


\bibitem{12} 


Бибербах Л. {\em Аналитическое продолжение}. -- М.: Наука, 1967. -- 239 с. 


\bibitem{13} 


Аксентьева Е.П., Гарифьянов Ф.Н. {\em К исследованию интегрального 


уравнения с ядром Карлемана} // Изв. вузов. Математика. -- 


1983. -- \No\,4. -- С.\,43--51. 


 


\end{thebibliography} 


 


\vskip 2cm 


 


\post{\it Казанский государственный}{\it Поступилa } 


\post{\it энергетический университет}{10.11.2002}


 


\label{end} 


\end{document} 


 


 


 


 


first version - in Jan 1997 !! 


 


